% Imporditud: tools/import_eksperimendid.py
% Allikas: Eesti füüsikaolümpiaadi eksperimendi ülesannete
%   kogu (Rael Kalda, 2023), ülesanne Ü49, lahendus L49.
\setAuthor{EFO žürii}
\setRound{piirkonnavoor}
\setYear{2007}
\setNumber{G E1}
\setDifficulty{2}
\setTopic{vedelike mehaanika}
\setVahendid{suur laia kaelaga purk, kitsas katseklaas, joonlaud, vesi}
\setPunktid{8}
\setAllikas{Eksperimendi kogu Ü49, lahendus lk 55}
\setLahendusseis{toores}

\prob{Katsesklaas}
Määrata katseklaasi mass.

\hint

\solu
Katseklaasi tuleb panna natuke vett, et see ujuks vertikaalasendis ja mõõta sisemise ning välimise veenivoo vahe $\Delta$h ning katseklaasi diameeter d. Siis võrdub katseklaasile mõjuv raskusjõud talle mõjuva üleslükkejõuga. Arvestades, et katseklaasis oleva vee panus on mõlemasse ühesugune, saame

mkg = $\rho$vV g $\Rightarrow$ mk = $\rho$vV = $\pi\rho$vd2$\Delta$h .

Siin mk on katseklaasi mass ja $\rho$v --- vee tihedus.

Hindamine: Idee --- 2 p, katseklaasi diameetri mõõtmine --- 1 p, ühekordne katse läbiviimine ja nivoode vahe mõõtmine --- 1 p, arvutused --- 2 p, korduvad katsed (vähemalt 3 korda) --- 1 p, piisavalt täpne lõppvastus --- 1 p.

Võimalik vale lahenduskäik: Paneme anumasse vett ja mõõdame veenivoo $h_0$ selles. Asetame katseklaasi anumasse ilma katseklaasi vett lisamata. Mõõdame anuma diameetri $d_0$ ja uue veenivoo anumas $h_1$. Siis oletades, et katseklaasi raskusjõu ja üleslükkejõu tasakaalu tõttu võrdub katseklaasi mass väljatõrjutud vee massiga, saame

mk = $\rho$vV = $\pi\rho$vd20 ($h_1$ $-$ $h_0$) .

Tegelikkuses, kui katses katseklaasi vett mitte panna, kukub ta pikali. Lisaks raskusjõule ja üleslükkejõule mõjub talle siis ka hõõrdejõud anuma seina poolt. Sel juhul aga massi määramiseks arvutusi teha pole enam võimalik, mistõttu sellist lahendust ei saa lugeda õigeks.

Hindamine: (kokku 4 p) Idee --- 1 p, katse teostamine ja mõõtmised --- 1 p, arvutused --- 1 p, korduvad katsed (vähemalt 3 korda) --- 1 p.
\probend
